% ===== 原稿 PDF 第 11 页 =====
% 首行为上页性质 7 的应用（证明收尾）；本页依次出现：性质 6、性质 7、性质 4（原稿编号次序即为 6、7、4）

由性质 7 $\Rightarrow$
\begin{gather*}
|A|=|C|=(-1)^{s+1}a_{s1}Q_{s1}\qquad\leftarrow\ \text{有 1 列全为 } 0\\
=0
\end{gather*}

\begin{proposition}[性质 6]
\[
\resizebox{\linewidth}{!}{$\displaystyle
\underbrace{\begin{vmatrix}
a_{11} & \cdots & a_{1r}+b_{1r} & \cdots & a_{1n}\\
a_{21} & \cdots & a_{2r}+b_{2r} & \cdots & a_{2n}\\
\vdots & & \vdots & & \vdots\\
a_{n1} & \cdots & a_{nr}+b_{nr} & \cdots & a_{nn}
\end{vmatrix}}_{|C|,\ Q_{ij}}
=
\underbrace{\begin{vmatrix}
a_{11} & \cdots & a_{1r} & \cdots & a_{1n}\\
a_{21} & \cdots & a_{2r} & \cdots & a_{2n}\\
\vdots & & \vdots & & \vdots\\
a_{n1} & \cdots & a_{nr} & \cdots & a_{nn}
\end{vmatrix}}_{|A|,\ M_{ij}}
+
\underbrace{\begin{vmatrix}
a_{11} & \cdots & b_{1r} & \cdots & a_{1n}\\
a_{21} & \cdots & b_{2r} & \cdots & a_{2n}\\
\vdots & & \vdots & & \vdots\\
a_{n1} & \cdots & b_{nr} & \cdots & a_{nn}
\end{vmatrix}}_{|B|,\ N_{ij}}$}
\]
\end{proposition}

\begin{proof}
$n=1$ 显然成立；设阶 $=n-1$ 时结论已成立，现证 $n$ 阶情形。

\noindent\textbf{Case 1}\quad $r=1$：
\[
|C|=(a_{11}+b_{11})Q_{11}-(a_{21}+b_{21})Q_{21}+\cdots+(-1)^{n+1}(a_{n1}+b_{n1})Q_{n1}
\]
注意 $Q_{i1}=M_{i1}=N_{i1},\ \forall i\ \Rightarrow\ |C|=|A|+|B|$。

\noindent\textbf{Case 2}\quad $r>1$：
\[
|C|=a_{11}Q_{11}-a_{21}Q_{21}+\cdots+(-1)^{n+1}a_{n1}Q_{n1}
\]
$Q_{i1},\,M_{i1},\,N_{i1}$ 满足性质 6 的条件 $\Rightarrow\ Q_{i1}=M_{i1}+N_{i1},\ \forall i$，
\[
\Rightarrow\ |C|=|A|+|B|.
\]
\end{proof}

\begin{proposition}[性质 7]
\[
\resizebox{\linewidth}{!}{$\displaystyle
\begin{vmatrix}
a_{11} & \cdots & a_{1r} & \cdots & a_{1s}+ca_{1r} & \cdots & a_{1n}\\
a_{21} & \cdots & a_{2r} & \cdots & a_{2s}+ca_{2r} & \cdots & a_{2n}\\
\vdots & & \vdots & & \vdots & & \vdots\\
a_{n1} & \cdots & a_{nr} & \cdots & a_{ns}+ca_{nr} & \cdots & a_{nn}
\end{vmatrix}
=
\begin{vmatrix}
a_{11} & \cdots & a_{1r} & \cdots & a_{1s} & \cdots & a_{1n}\\
a_{21} & \cdots & a_{2r} & \cdots & a_{2s} & \cdots & a_{2n}\\
\vdots & & \vdots & & \vdots & & \vdots\\
a_{n1} & \cdots & a_{nr} & \cdots & a_{ns} & \cdots & a_{nn}
\end{vmatrix}
+
\underbrace{\begin{vmatrix}
a_{11} & \cdots & a_{1r} & \cdots & ca_{1r} & \cdots & a_{1n}\\
a_{21} & \cdots & a_{2r} & \cdots & ca_{2r} & \cdots & a_{2n}\\
\vdots & & \vdots & & \vdots & & \vdots\\
a_{n1} & \cdots & a_{nr} & \cdots & ca_{nr} & \cdots & a_{nn}
\end{vmatrix}}_{0}$}
\]
\end{proposition}

\begin{proof}
性质 6 $+$ 性质 3。
\end{proof}

\begin{proposition}[性质 4]
对换 $|A|$ 的两个不同列，得到行列式 $|B|=-|A|$。
\end{proposition}

\begin{proof}
设 $|B|$ 是对换 $|A|$ 的第 $r$ 列与第 $s$ 列所得：
\[
\resizebox{\linewidth}{!}{$\displaystyle
|A|=\begin{vmatrix}
a_{11} & \cdots & a_{1r} & \cdots & a_{1s} & \cdots & a_{1n}\\
a_{21} & \cdots & a_{2r} & \cdots & a_{2s} & \cdots & a_{2n}\\
\vdots & & \vdots & & \vdots & & \vdots\\
a_{n1} & \cdots & a_{nr} & \cdots & a_{ns} & \cdots & a_{nn}
\end{vmatrix}
\ \longrightarrow\
|B|=\begin{vmatrix}
a_{11} & \cdots & a_{1s} & \cdots & a_{1r} & \cdots & a_{1n}\\
a_{21} & \cdots & a_{2s} & \cdots & a_{2r} & \cdots & a_{2n}\\
\vdots & & \vdots & & \vdots & & \vdots\\
a_{n1} & \cdots & a_{ns} & \cdots & a_{nr} & \cdots & a_{nn}
\end{vmatrix}$}
\]
% ===== 原稿 PDF 第 12 页 =====
构造一个新行列式
\[
|C|=\begin{vmatrix}
a_{11} & \cdots & a_{1r}+a_{1s} & \cdots & a_{1r}+a_{1s} & \cdots & a_{1n}\\
a_{21} & \cdots & a_{2r}+a_{2s} & \cdots & a_{2r}+a_{2s} & \cdots & a_{2n}\\
\vdots & & \vdots & & \vdots & & \vdots\\
a_{n1} & \cdots & a_{nr}+a_{ns} & \cdots & a_{nr}+a_{ns} & \cdots & a_{nn}
\end{vmatrix}
\]
由性质 6，
\[
\resizebox{\linewidth}{!}{$\displaystyle
|C|=
\begin{vmatrix}
a_{11} & \cdots & a_{1r} & \cdots & a_{1r} & \cdots & a_{1n}\\
a_{21} & \cdots & a_{2r} & \cdots & a_{2r} & \cdots & a_{2n}\\
\vdots & & \vdots & & \vdots & & \vdots\\
a_{n1} & \cdots & a_{nr} & \cdots & a_{nr} & \cdots & a_{nn}
\end{vmatrix}
+|A|+|B|+
\begin{vmatrix}
a_{11} & \cdots & a_{1s} & \cdots & a_{1s} & \cdots & a_{1n}\\
a_{21} & \cdots & a_{2s} & \cdots & a_{2s} & \cdots & a_{2n}\\
\vdots & & \vdots & & \vdots & & \vdots\\
a_{n1} & \cdots & a_{ns} & \cdots & a_{ns} & \cdots & a_{nn}
\end{vmatrix}$}
\]
% 原稿上式最末一个行列式的末列作 a_{ir}（即写成 |B|），该行列式两列相同应为 0，此处按 a_{is} 记
\[
\Rightarrow\ |B|=-|A|.
\]
\end{proof}

考虑如下相邻对换：
\[
1\ \ 2\ \ \cdots\ \ r-1\ \ r\ \ \cdots\ \ n
\qquad\Rightarrow\qquad
r\ \ 1\ \ 2\ \ \cdots\ \ r-1\ \ r+1\ \ \cdots\ \ n
\]
\[
\underbrace{|A|}_{M_{ij}}
\ \xrightarrow{\ r-1\ \text{次列对换}\ }\
\underbrace{|B|}_{N_{ij}}
\qquad\Rightarrow\qquad
|B|=(-1)^{r-1}|A|
\]
\[
\resizebox{\linewidth}{!}{$\displaystyle
|B|=\begin{vmatrix}
a_{1r} & a_{11} & \cdots & a_{1,r-1} & a_{1,r+1} & \cdots & a_{1n}\\
a_{2r} & a_{21} & \cdots & a_{2,r-1} & a_{2,r+1} & \cdots & a_{2n}\\
\vdots & \vdots & & \vdots & \vdots & & \vdots\\
a_{nr} & a_{n1} & \cdots & a_{n,r-1} & a_{n,r+1} & \cdots & a_{nn}
\end{vmatrix}$}
\]
\[
\Rightarrow\ |B|=a_{1r}N_{11}-a_{2r}N_{21}+\cdots+(-1)^{n+1}a_{nr}N_{n1}
\]
\[
\forall i,\qquad N_{i1}=M_{ir}
\]
\begin{align*}
\Rightarrow\ |A|&=(-1)^{r+1}|B|=(-1)^{r-1}\Bigl(a_{1r}M_{1r}-a_{2r}M_{2r}+\cdots+(-1)^{n+1}a_{nr}M_{nr}\Bigr)\\
&=(-1)^{r+1}a_{1r}M_{1r}+(-1)^{r+2}a_{2r}M_{2r}+\cdots+(-1)^{r+n}a_{nr}M_{nr}
\end{align*}
\[
\boxed{\;|A|=a_{1r}A_{1r}+a_{2r}A_{2r}+\cdots+a_{nr}A_{nr}\;}
\qquad\text{按第 } r \text{ 列展开}.
\]

Kronecker 符号
\[
\delta_{ij}=
\begin{cases}
1, & \text{若 } i=j,\\
0, & \text{若 } i\neq j.
\end{cases}
\]

\begin{theorem}% 原稿标为「定理 1」
设 $|A|$ 为 $n$ 阶行列式，$1\le r,\,s\le n$，则
\[
a_{1r}A_{1s}+a_{2r}A_{2s}+\cdots+a_{nr}A_{ns}=\delta_{rs}|A|.
\]
\end{theorem}

% ===== 原稿 PDF 第 13 页 =====
\begin{proof}
若 $r=s$，已证；下不妨 $r<s$。

构造一个新的行列式
\[
0=|C|=\begin{vmatrix}
a_{11} & \cdots & a_{1r} & \cdots & a_{1r} & \cdots & a_{1n}\\
a_{21} & \cdots & a_{2r} & \cdots & a_{2r} & \cdots & a_{2n}\\
\vdots & & \vdots & & \vdots & & \vdots\\
a_{n1} & \cdots & a_{nr} & \cdots & a_{nr} & \cdots & a_{nn}
\end{vmatrix}
=a_{1r}A_{1s}+a_{2r}A_{2s}+\cdots+a_{nr}A_{ns}.
\]
（$|C|$ 按第 $s$ 列展开）
\end{proof}

\begin{lemma}\label{lem:2}% 原稿标为「引理 2」
\[
|A|=\begin{vmatrix}
a_{11} & \cdots & a_{1r} & \cdots & a_{1n}\\
\vdots & & \vdots & & \vdots\\
0 & \cdots & a_{sr} & \cdots & 0\\
\vdots & & \vdots & & \vdots\\
a_{n1} & \cdots & a_{nr} & \cdots & a_{nn}
\end{vmatrix}
=a_{sr}A_{sr}.
\]
\end{lemma}

\begin{proof}
$|A|$ 按第 $r$ 列进行展开：
\[
|A|=a_{1r}A_{1r}+\cdots+a_{sr}A_{sr}+\cdots+a_{nr}A_{nr}.
\]
$\forall t\neq s$，$A_{tr}=0$（$A_{tr}$ 中至少有一行全为 $0$），
\[
\Rightarrow\ |A|=a_{sr}A_{sr}.
\]
\end{proof}

\begin{lemma}% 原稿标为「引理 3」
\[
|A|=a_{r1}A_{r1}+a_{r2}A_{r2}+\cdots+a_{rn}A_{rn}
\]
按任何一行进行展开。
\end{lemma}

\begin{proof}
第 $r$ 行元素的拆分：
\[
\begin{aligned}
a_{r1}&=a_{r1}+0+\cdots+0,\\
a_{r2}&=0+a_{r2}+\cdots+0,\\
a_{rn}&=0+0+\cdots+a_{rn}.
\end{aligned}
\]
\[
\resizebox{\linewidth}{!}{$\displaystyle
\Rightarrow\ |A|=
\begin{vmatrix}
a_{11} & a_{12} & \cdots & a_{1n}\\
\vdots & \vdots & & \vdots\\
a_{r1} & 0 & \cdots & 0\\
\vdots & \vdots & & \vdots\\
a_{n1} & a_{n2} & \cdots & a_{nn}
\end{vmatrix}
+
\begin{vmatrix}
a_{11} & a_{12} & \cdots & a_{1n}\\
\vdots & \vdots & & \vdots\\
0 & a_{r2} & \cdots & 0\\
\vdots & \vdots & & \vdots\\
a_{n1} & a_{n2} & \cdots & a_{nn}
\end{vmatrix}
+\cdots+
\begin{vmatrix}
a_{11} & a_{12} & \cdots & a_{1n}\\
\vdots & \vdots & & \vdots\\
0 & 0 & \cdots & a_{rn}\\
\vdots & \vdots & & \vdots\\
a_{n1} & a_{n2} & \cdots & a_{nn}
\end{vmatrix}$}
\]
由引理~\ref{lem:2}，
\[
=a_{r1}A_{r1}+a_{r2}A_{r2}+\cdots+a_{rn}A_{rn}.
\]
\end{proof}

% ===== 原稿 PDF 第 14 页 =====
\begin{theorem}% 原稿标为「定理 4」
设 $|A|$ 为 $n$ 阶行列式，$1\le r,\,s\le n$，则
\[
a_{r1}A_{s1}+a_{r2}A_{s2}+\cdots+a_{rn}A_{sn}=\delta_{rs}|A|.
\]
\end{theorem}

\begin{proof}
$r=s$ 已证，下不妨设 $r<s$。

构造新行列式
\[
0=|C|=\begin{vmatrix}
a_{11} & a_{12} & \cdots & a_{1n}\\
\vdots & \vdots & & \vdots\\
a_{r1} & a_{r2} & \cdots & a_{rn}\\
\cdots & \cdots & \cdots & \cdots\\
a_{r1} & a_{r2} & \cdots & a_{rn}\\
a_{n1} & a_{n2} & \cdots & a_{nn}
\end{vmatrix}
=a_{r1}A_{s1}+a_{r2}A_{s2}+\cdots+a_{rn}A_{sn}.
\]
（$|C|$ 按第 $s$ 行展开）
\end{proof}

\begin{definition}[转置行列式]% 原稿标为「定义 5」
\[
\resizebox{\linewidth}{!}{$\displaystyle
|A|=\underbrace{\begin{vmatrix}
a_{11} & a_{12} & \cdots & a_{1n}\\
a_{21} & a_{22} & \cdots & a_{2n}\\
\vdots & \vdots & & \vdots\\
a_{n1} & a_{n2} & \cdots & a_{nn}
\end{vmatrix}}_{M_{ij}\text{（余子式）}}
\qquad
|A'|=\underbrace{\begin{vmatrix}
a_{11} & a_{21} & \cdots & a_{n1}\\
a_{12} & a_{22} & \cdots & a_{n2}\\
\vdots & \vdots & & \vdots\\
a_{1n} & a_{2n} & \cdots & a_{nn}
\end{vmatrix}}_{N_{ij}\text{（余子式）}}$}
\]
称为 $A$ 的转置。即 $|A|$ 的第 $i$ 行就是 $|A'|$ 的第 $i$ 列（$\forall i$）。
\end{definition}

\begin{proposition}[性质 8]
\[
|A'|=|A|.
\]
\end{proposition}

\begin{proof}
对阶数进行归纳：$n=1$ $\checkmark$；设 $n-1$ 阶 $\checkmark$，下证 $n$ 阶。
\[
|A|=a_{11}M_{11}-a_{21}M_{21}+\cdots+(-1)^{n+1}a_{n1}M_{n1}.
\]
% 原稿此处作「N_{ij} 是 M_{ij} 的转置」，按后文 N_{ji}=M_{ij}，应为 N_{ji} 是 M_{ij} 的转置
$\forall i,j$，$N_{ji}$ 是 $M_{ij}$ 的转置，由归纳假设 $N_{ji}=M_{ij}$，$\forall i,j$。
\[
|A|=a_{11}N_{11}-a_{21}N_{12}+\cdots+(-1)^{n+1}a_{n1}N_{1n}=|A'|
\]
（按第 $1$ 行进行展开）。
\end{proof}

讨论 $n$ 元线性方程组
\[
(*)\quad
\begin{cases}
a_{11}x_1+a_{12}x_2+\cdots+a_{1n}x_n=b_1,\\
a_{21}x_1+a_{22}x_2+\cdots+a_{2n}x_n=b_2,\\
\qquad\cdots\cdots\\
a_{n1}x_1+a_{n2}x_2+\cdots+a_{nn}x_n=b_n,
\end{cases}
\]
设 $(*)$ 有解。

\[
\text{系数行列式}\qquad
|A|=\begin{vmatrix}
a_{11} & a_{12} & \cdots & a_{1n}\\
a_{21} & a_{22} & \cdots & a_{2n}\\
\vdots & \vdots & & \vdots\\
a_{n1} & a_{n2} & \cdots & a_{nn}
\end{vmatrix}.
\]
\[
\resizebox{\linewidth}{!}{$\displaystyle
|A_1|=\begin{vmatrix}
b_1 & a_{12} & \cdots & a_{1n}\\
b_2 & a_{22} & \cdots & a_{2n}\\
\vdots & \vdots & & \vdots\\
b_n & a_{n2} & \cdots & a_{nn}
\end{vmatrix}
\overset{\text{替换}}{=\joinrel=}
\begin{vmatrix}
a_{11}x_1+a_{12}x_2+\cdots+a_{1n}x_n & a_{12} & \cdots & a_{1n}\\
a_{21}x_1+a_{22}x_2+\cdots+a_{2n}x_n & a_{22} & \cdots & a_{2n}\\
\vdots & \vdots & & \vdots\\
a_{n1}x_1+a_{n2}x_2+\cdots+a_{nn}x_n & a_{n2} & \cdots & a_{nn}
\end{vmatrix}
=
\begin{vmatrix}
a_{11}x_1 & a_{12} & \cdots & a_{1n}\\
a_{21}x_1 & a_{22} & \cdots & a_{2n}\\
\vdots & \vdots & & \vdots\\
a_{n1}x_1 & a_{n2} & \cdots & a_{nn}
\end{vmatrix}
=x_1|A|$}
\]
若 $(*)$ 有解，
\[
\Rightarrow\quad x_1=\frac{|A_1|}{|A|},\ \cdots,\ x_n=\frac{|A_n|}{|A|}.
\]

% ===== 原稿 PDF 第 15 页 =====
\begin{theorem}% 原稿标为「定理 6」
（Cramer 法则）% 原稿写作「Cramon 法则」
若 $|A|\neq 0$，则 $(*)$ 有唯一解
\[
x_i=\frac{|A_i|}{|A|}\qquad (i=1,\dots,n).
\]
\end{theorem}

\begin{proof}
上面告诉：若 $(*)$ 有解，则解必为如上形式。只要证明 $x_i=\dfrac{|A_i|}{|A|}$ 确为 $(*)$ 的解即可。\renewcommand{\qedsymbol}{}
\end{proof}

\section{行列式的计算（上）}% 原稿节标题块作「7-8 行列式的计算(上) / 行列式的计算(下)」

\[
(*)\quad
\begin{cases}
a_{11}x_1+a_{12}x_2+\cdots+a_{1n}x_n=b_1,\\
a_{21}x_1+a_{22}x_2+\cdots+a_{2n}x_n=b_2,\\
\qquad\cdots\cdots\\
a_{n1}x_1+a_{n2}x_2+\cdots+a_{nn}x_n=b_n,
\end{cases}
\]

\textbf{Cramer 法则}\quad 若 $|A|\neq 0$，则 $(*)$ 有且仅有一组解
\[
x_1=\frac{|A_1|}{|A|},\ \cdots,\ x_n=\frac{|A_n|}{|A|},
\]
其中
\[
|A_j|=\begin{vmatrix}
a_{11} & \cdots & b_1 & \cdots & a_{1n}\\
a_{21} & \cdots & b_2 & \cdots & a_{2n}\\
\vdots & & \vdots & & \vdots\\
a_{n1} & \cdots & b_n & \cdots & a_{nn}
\end{vmatrix}.
\]

\begin{proof}
只要验证 $x_j=\dfrac{|A_j|}{|A|}$ 确为 $(*)$ 的解。
\[
a_{r1}A_{s1}+a_{r2}A_{s2}+\cdots+a_{rn}A_{sn}=\delta_{rs}|A|
\]
\[
\sum_{j=1}^{n}a_{rj}A_{sj}=\delta_{rs}|A|
\]
\[
\resizebox{\linewidth}{!}{$\displaystyle
\begin{array}{r@{\quad}c}
\sum\limits_{j=1}^{n}a_{1j} & a_{11}\quad a_{12}\quad\cdots\quad a_{1n}\\[4pt]
+\,\sum\limits_{j=1}^{n}a_{2j} & a_{21}\quad a_{22}\quad\cdots\quad a_{2n}\\[4pt]
+\,\vdots & \vdots\\[4pt]
+\,\sum\limits_{j=1}^{n}a_{mj} & a_{m1}\quad a_{m2}\quad\cdots\quad a_{mn}\\[4pt]
=\,\sum\limits_{i=1}^{m}\sum\limits_{j=1}^{n}a_{ij} &
\Bigl(\sum\limits_{i=1}^{m}a_{i1}\quad \sum\limits_{i=1}^{m}a_{i2}\quad\cdots\quad \sum\limits_{i=1}^{m}a_{in}\Bigr)=\sum\limits_{j=1}^{n}\sum\limits_{i=1}^{m}a_{ij}
\end{array}$}
\]
\[
\sum_{j=1}^{n}\sum_{i=1}^{m}a_{ij}=\sum_{i=1}^{m}\sum_{j=1}^{n}a_{ij}.
\]
\[
x_j=\frac{|A_j|}{|A|}=\frac{1}{|A|}\sum_{i=1}^{n}b_iA_{ij}
\]
\renewcommand{\qedsymbol}{}% 证明延续至下页，此处不打印证毕符号
\end{proof}
